\section{Method}

\subsection{The Ensemble Puzzle}

Standard practice for improving accuracy is
ensembling~\citep{dietterich2000ensemble}: train multiple models with different
seeds and average their predictions. For TRM on Sudoku-Extreme:
%
% Source: x02.log (86.02%), x03.log (83.87%), x04.log (85.03%) - noised models
\begin{center}
\begin{tabular}{lcc}
\toprule
Method & Puzzle Acc. & Forward Passes \\
\midrule
Single model & 84.97\% & 16 \\  % avg(x02.log, x03.log, x04.log)
Ensemble (3 seeds) & 88.52\% & 48 \\  % x02.log + x03.log + x04.log
+ TTA (4 rotations) & 91.51\% & 192 \\  % x02.log + x03.log + x04.log + TTA
\bottomrule
\end{tabular}
\end{center}
%
Ensemble plus test-time augmentation yields +6.5\% accuracy, but at
12$\times$ the compute cost. This presents a puzzle: ensembling helps, but the
standard approach of probability averaging seems wasteful. Each model in the
ensemble has its own ``opinion'' about when to stop reasoning, yet we ignore
this timing information entirely. Can we do better by paying attention to
\emph{when} models commit to their answers?

We begin with a biological intuition and follow it to its logical conclusion.

Biological neural systems face a fundamental constraint: metabolic energy is
scarce, yet decisions must be fast. Evolution's solution is to \emph{act on
the first confident signal}. In cortical winner-take-all
circuits~\citep{grossberg1987competitive}, competing neural populations race to
threshold, with the first to fire suppressing alternatives. Spiking neural
networks formalize this as time-to-first-spike (TTFS)
coding~\citep{thorpe2001spike, park2020t2fsnn}, where information is encoded in
\emph{when} neurons fire rather than \emph{how often}---achieving both speed
and energy efficiency.

This principle suggests a strategy for ensembles of iterative reasoners:
instead of averaging probabilities, \textbf{select the first model to halt}.

\subsection{Halt-First Ensemble}

The procedure is simple: run all ensemble members in parallel, each performing
its own ACT iterations. When any model's halting signal exceeds threshold, that
model provides the final answer. The results are striking:
%
% Source: x02.log, x03.log, x04.log (3 seeds × 4 rotations = 12 chains)
\begin{center}
\begin{tabular}{@{}lccc@{}}
\toprule
Method & Acc. & Steps & Speedup \\
\midrule
Prob.\ avg.\ (12) & 91.5\% & 192 & 1$\times$ \\  % x02.log + x03.log + x04.log + TTA
\textbf{Halt-first (12)} & \textbf{97.2\%} & \textbf{18.5} & \textbf{10$\times$} \\  % x02.log + x03.log + x04.log + TTA halt-first
\bottomrule
\end{tabular}
\end{center}
%
Halt-first selection improves accuracy by 5.7 percentage points while using
10$\times$ less compute. This validates the biological intuition:
\textbf{inference speed is an implicit confidence signal}. When a reasoning
process converges quickly, it indicates the model has found a ``clean''
solution path---one without contradictions or backtracking. Probability
averaging destroys this signal; halt-first selection preserves the information
that evolution discovered: the reasoner that finishes first is most likely
correct. We analyze the halting dynamics in detail in \Cref{sec:halting}.

\subsection{Oracle-First Training}

The halt-first result raises a natural question: can we internalize this
benefit within a single model? Deploying multiple models is expensive---we want
the accuracy of an ensemble with the deployment cost of one.

Our key insight is that ensemble diversity comes from different random
initializations during training. We can internalize this diversity by
maintaining $K$ parallel latent initializations \emph{within} one model and
letting them compete. During training, we run $K{=}4$ parallel hypotheses; at
inference, we run just one.

Specifically, we maintain $K$ parallel low-level states
$\{z_L^{(k)}\}_{k=1}^K$, each from a learned initialization
$L_{\text{init}}^{(k)}$, while sharing the high-level state $z_H$. At each
step, all $K$ heads process the input in parallel, and the
\textbf{winner-take-all} (WTA) rule selects which head receives the gradient:
%
\begin{equation}
  k^* = \argmin_k \mathcal{L}(y^{(k)}, y_\text{target})
\end{equation}
%
Only the winning head contributes to the loss. This mirrors halt-first
selection: at inference, the fastest model wins; during training, the most
accurate head wins. Both reward the same underlying property: having found a
clean solution path. The connection is deep---we are teaching the model to find
confident solutions by making confidence (low loss) the criterion for receiving
gradient signal.

\paragraph{Training dynamics.} A key detail: TRM training ``unpacks'' the outer
loop of \Cref{alg:trm}. Each gradient step performs one H-cycle, carrying
$z_H$ and $z_L$ to the next step. This means the $K$ parallel chains evolve
across multiple gradient updates, not within a single forward pass.

\paragraph{Carry policy.} After each step, we must decide how to propagate
states to the next iteration. We use \texttt{copy\_top1} for $z_H$: the
winner's $z_H$ is copied to all chains (\Cref{alg:wta}, line 8). For $z_L$, we
use \texttt{all}: each chain keeps its own $z_L$. This allows heads to benefit
from shared high-level progress while maintaining low-level diversity (see
\Cref{sec:ablations} for ablations).

\begin{algorithm}[h]
\caption{Oracle-First Training (one step)}
\label{alg:wta}
\begin{algorithmic}[1]
  \STATE \textbf{Input:} $x$, $y_\text{target}$, carried $z_H$, $\{z_L^{(k)}\}_{k=1}^K$
  \FOR{$k = 1$ to $K$}
    \STATE $y_k, q_k, z_H', z_L^{(k)} \gets \operatorname{TRM}(x, z_H, z_L^{(k)}; \theta)$
  \ENDFOR
  \STATE $k^* \gets \argmin_{k} \mathcal{L}_\text{CE}(y_k, y_\text{target})$ \LCOMMENT{winner}
  \STATE $c \gets \mathbf{1}[y_{k^*} = y_\text{target}]$ \LCOMMENT{correctness}
  \STATE $\mathcal{L} \gets \mathcal{L}_\text{CE}(y_{k^*}, y_\text{target}) + \lambda \mathcal{L}_\text{BCE}(\sigma(q_{k^*}), c)$
  \STATE Backprop through $\mathcal{L}$
  \STATE $z_H \gets z_H'$ \LCOMMENT{winner's $z_H$ shared to all}
\end{algorithmic}
\end{algorithm}
%
\noindent Here $\mathcal{L}_\text{CE}$ is cross-entropy with label smoothing
($\alpha{=}0.2$), $\mathcal{L}_\text{BCE}$ is binary cross-entropy for the
halting signal where $q_k$ is the halting logit and $\sigma(q_k)$ is the
halting probability, and $\lambda{=}0.05$.

\paragraph{Counter-intuitive design.} WTA training deliberately spends compute
on reasoning chains that won't contribute to the gradient. This apparent waste
enables exploration: losing heads probe alternative paths, and winner selection
identifies which succeeded.

\subsection{Faster Training}

WTA training requires $K$ forward passes per iteration---a $4\times$ increase
for $K{=}4$. With only a single GPU, efficient training is essential. We
combine two techniques that reduce training time from days to hours, making
rapid experimentation possible.

\subsubsection{Muon Optimizer}

We use Muon~\citep{jordan2024muon} for 2D weight matrices (lr${=}$0.02,
wd${=}$0.005) and AdamW for embeddings, heads, and biases (lr${=}10^{-4}$,
wd${=}$1.0, $\beta{=}$(0.9, 0.95)). Muon's orthogonalized momentum enables
substantially higher learning rates than AdamW alone, accelerating convergence.
However, Muon is magnitude-invariant, so bias-like parameters require AdamW.
Learning rate follows a cosine schedule.

\paragraph{Selective regularization.} This split-optimizer design creates an
asymmetric regularization structure: the core reasoning weights (MLPMixer
blocks) receive minimal regularization via Muon's low weight decay, while
embeddings and output heads receive strong regularization via AdamW's high
weight decay. Combined with label smoothing ($\alpha{=}0.2$), this encourages
the core model to learn expressive representations while preventing the
input/output interfaces from overfitting. In contrast, baseline TRM applies
uniform high weight decay (1.0) to all parameters, which may over-constrain
the reasoning layers.

\paragraph{Modified SwiGLU.} We discovered that Muon fails completely with
standard SwiGLU activation. The fix is to apply layer normalization before the
down-projection: $\text{SwiGLU}_\text{mod}(x) = \text{SiLU}(g) \cdot
\text{LayerNorm}(g \odot x)$ where $g = Wx$. Without this modification, Muon's
orthogonalized updates cause training to diverge.

\subsubsection{SVD-Aligned Initialization}

The $K$ initialization vectors $L_{\text{init}}^{(k)}$ determine hypothesis
diversity. Naive random initialization risks ``dead'' heads---initializations
that point in directions the network cannot use effectively. We address this by
computing the top singular vectors of the first layer's weight matrix and
ensuring all $K$ heads have equal alignment with this subspace. This provides
equal ``effective signal strength'' across heads while maintaining diversity in
the orthogonal complement.

\subsection{Intuition: Why Speed Indicates Confidence}

We offer a speculative interpretation connecting halting speed to optimization
conditioning. This is a hypothesis, not a formal result; empirical validation
is future work.

\paragraph{Bilevel fixed-point structure.} TRM computes a nested equilibrium:
%
\begin{align}
  z_L^* &= f(z_L^*, z_H, z_x) && \text{(inner: fast local refinement)} \\
  z_H^* &= g(z_H^*, z_L^*) && \text{(outer: slow global integration)}
\end{align}
%
The inner loop equilibrates $z_L$ for fixed $z_H$; the outer loop updates $z_H$
given the equilibrated $z_L^*$. This separation creates a natural hierarchy:
$z_L$ handles cell-level constraint propagation while $z_H$ maintains
puzzle-level state.

\paragraph{Architectural asymmetry stabilizes the inner loop.} A key
observation: $z_x$, $z_L$, and $z_H$ each have unit variance (due to
normalization), but they enter the computation asymmetrically. L-cycles receive
$z_L + z_H + z_x$ (three terms), while H-cycles receive only $z_H + z_L$ (two
terms). Assuming approximate independence, $z_L$ contributes $\tfrac{1}{3}$ of
L-cycle input variance versus $\tfrac{1}{2}$ for H-cycles.

The Jacobian $J_L = \partial f / \partial z_L$ measures sensitivity to $z_L$
perturbations. When $z_L$ is a smaller fraction of the input, the output is
more ``anchored'' by context $(z_H + z_x)$ and less sensitive to $z_L$---yielding
smaller spectral radius. The fixed input $z_x$ thus acts as a stabilizing
anchor that L-cycles have but H-cycles lack, naturally making the inner loop
better conditioned than the outer loop.

\paragraph{WTA training enhances this effect.} The \texttt{copy\_top1} carry
policy (\Cref{alg:wta}, line~8) copies the winner's $z_H$ to all chains. This
means each chain's $z_H$ comes from a \emph{different} chain than its $z_L$,
further decorrelating the inputs and strengthening the independence assumption
underlying the variance argument.

\paragraph{Speed reflects conditioning (conjecture).} For iterative fixed-point
methods, convergence rate depends on the spectral radius of the iteration
Jacobian---smaller spectral radius means faster convergence. We hypothesize
that the learned halting signal $q_\text{halt}$ correlates with this
conditioning: chains that halt quickly have found well-conditioned paths, while
chains that deliberate longer are stuck in poorly-conditioned regions.
Evolution's heuristic (trust the first signal) may thus correspond to selecting
well-conditioned solutions.
